\section{A square with a small marked moment}
\label{sec:weights}

We construct the weight in Proposition~\ref{prop:weights} as the square
of a signed sum of smooth divisor sums. We first define that square
and state the exact analytic estimates it satisfies. The next section
uses their formulas to choose the coefficients and functions; the
proofs of the analytic estimates follow in Section~\ref{sec:moments}.

Fix $\lambda>0$, and write
\[
 L=\log X,\qquad h=\lfloor\lambda L\rfloor,\qquad
 J=\{1,\ldots,h\},\qquad I=\{h+1,\ldots,2h\}.
\]
Here $X$ tends to infinity through integers, and $\mathbb E_X$ denotes
the average over the integers $X<m\le2X$.  Recall that
\[
 V_B(m)=\frac1L\sum_{b\in B}\vartheta(m+b),\qquad
 \vartheta(n)=(\log n)1_{\{n\ {\text{prime}}\}}.
\]
Let $\mu$ denote the M\"obius function. For a compactly supported
function $F:[0,\infty)^j\to\mathbb R$ and a tuple of distinct shifts
$\mathbf b=(b_1,\ldots,b_j)$, define
\[
 D_F(m;\mathbf b)
 =\sum_{d_i\mid m+b_i}\mu(d_1)\cdots\mu(d_j)
       F\left(\frac{\log d_1}{L},\ldots,
              \frac{\log d_j}{L}\right).
\]
The sum runs over positive divisors.
The zero-dimensional convention makes $D_F=F$ a scalar.  A symmetric
function $F$ allows us to write $D_F(m;S)$ for an unordered shift set $S$.

Put $\tau=1/8$, and let $k$ be a positive integer. For $1\le j\le k$, take a real symmetric function
\[
 f_j\in C_c^\infty
 \bigl(\{\mathbf t\in(0,\infty)^j:\textstyle\sum_i t_i<\tau\}\bigr).
\]
Allow also a real scalar $f_0$, and set $f_{k+1}=F_{k+1}=0$.  Define
\[
 F_j(\mathbf v)=\int_{s_i\ge v_i\ (1\le i\le j)}
                         f_j(\mathbf s)\,d\mathbf s,
 \qquad F_0=f_0,
 \qquad
 (Tf_{j+1})(\mathbf t)=\int_0^\infty f_{j+1}(\mathbf t,u)\,du.
\]
In dimension zero the last expression is a scalar.  All norms below
are Lebesgue $L^2$ norms on nonnegative orthants, with absolute value
in dimension zero. On the nonnegative orthant, $F_j$ vanishes when
$\sum_i v_i\ge\tau$. Cumulative integration also gives
\begin{equation}
 (-1)^j\partial_1\cdots\partial_jF_j=f_j.
 \label{eq:cumulative-derivative}
\end{equation}
Thus these complete mixed derivatives vanish on every coordinate
face, even though the cumulative functions themselves need not.

Define the signed sum
\[
 Z(m)=\sum_{j=0}^k\sum_{\substack{S\subset I\\|S|=j}}D_{F_j}(m;S),
 \qquad \alpha_j=\frac{\lambda^j}{j!}.
\]
The weight will be a positive normalization of $Z^2$. The following
analytic statement gives its mass, its interaction with the two prime
counts, and the second-moment bounds for prime detection and removal
of the weight. Its explicit formulas identify the remaining choice
of the functions $f_j$.

\begin{proposition}[Moment identities and a detector bound]
\label{prop:weight-moments}
Hold $\lambda$, $k$, and the family $(f_j)_{0\le j\le k}$ fixed.
Then, as $X\to\infty$,
\begin{align}
 \mathbb E_XZ^2&\longrightarrow
       w:=\sum_{j=0}^k\alpha_j\|f_j\|_2^2,
       &\mathbb E_XZ^4&=O_{\lambda,k,\mathbf f}(1),
       \label{eq:weight-unmarked}\\
 \mathbb E_X(Z^2V_J)&\longrightarrow\lambda w,
       &\mathbb E_X(Z^2V_I)&\longrightarrow
       v:=\lambda\sum_{j=0}^k\alpha_j
                         \|f_j+Tf_{j+1}\|_2^2.
       \label{eq:weight-marked}
\end{align}
Choose once and for all a real $G\in C_c^\infty(\mathbb R)$ with
$G(0)=1$ and $G(u)=0$ for $u\ge\tau$, and put
\[
 c_G=\int_0^\infty |G'(u)|^2\,du,
 \qquad C_*=\lambda+\lambda^2c_G^2.
\]
For every fixed family as above,
\begin{equation}
 \limsup_{X\to\infty}\mathbb E_X(Z^2V_J^2)\le C_*w.
 \label{eq:detector-bound}
\end{equation}
In particular, $C_*$ is independent of $k$ and of the family $(f_j)$.
\end{proposition}

The formulas distinguish two effects of a prime mark. A prime in $J$
lies outside every subset used to build $Z$, whereas a prime in $I$
alters the sum itself. To see the latter point exactly, fix $a\in I$.
When $m+a$ is prime, every nonzero term from a subset containing $a$
has divisor one in that coordinate: the other possible divisor satisfies
$\log(m+a)/L>1>\tau$, so its contribution vanishes by the support
condition on $F_j$.
Put $\widetilde F_j=F_j+F_{j+1}(\cdot,0)$. Deleting $a$ from each
subset that contains it therefore gives the exact identity
\begin{equation}\label{eq:prime-deletion}
 Z(m)=\sum_{j=0}^k
       \sum_{\substack{U\subset I\setminus\{a\}\\|U|=j}}
       D_{\widetilde F_j}(m;U)
 \qquad\text{when $m+a$ is prime}.
\end{equation}
Each subset occurs once, so deletion introduces no dimension factor.
The signed complete mixed derivative of $\widetilde F_j$ is
$f_j+Tf_{j+1}$. This explains why the marked form $v$ couples neighboring
dimensions.

The analytic proof evaluates the products in these squares by pairing
equal shifts. Any unmatched shift sets a coordinate of a complete
mixed derivative to zero, so the face vanishing in
\eqref{eq:cumulative-derivative} removes its main term. Only matching
subsets remain in the limiting quadratic forms, producing the squared
norms in $w$ and $v$. The same proof controls the errors and establishes
the two different second-moment bounds.

For a family with $w>0$, the normalization $W_X=Z^2/w$ therefore has
all the properties in Proposition~\ref{prop:weights} except the
prescribed smallness of the prime mass in $I$. That last requirement
is exactly $v/w\le\varepsilon$. The next section constructs families
with $w\to1$ and $v\to0$, while the detector constant $C_*$ remains
unchanged.
